% ============================================================
% Chapter 5: Experiments — 本文（Phase 2, 2026-08-09）
% 骨子版（Phase 1）は git 履歴の commit bf41ec4 以前を参照。
% 用語は §4.2 に合わせた（single-source / cross-source / synthesized）。
% 数値の一次ソースは各段落の % source: コメントに記載。
% 2026-10-21 改訂（日本語版 8/31 報告の 2 回目の改訂 4781cdd・再実験 c7c4feb にそろえた）:
%   (a) 設定 A は 91 件を除いた 1,732 件。拡張ケースは訓練にも評価にも使っていない（旧「拡張を入れると 0.521 → 0.622」の
%       段落は、拡張ケースを訓練にもテストにも入れた古い実験の値なので削除した）。設定 B は「正解式の数が 1〜10 本の
%       ケースが 100 件ずつ」と書いた（加藤先生の注釈 R23）。
%   (b) 定数を検証用のケースで選ぶ（baseline の重み、BM25・E5 の重み、提案手法の第 1 段の重みと不正解の選び方）。
%       第 1 段の重みは (0.3, 0.7) に固定し N_1 = 2,000（本人の決定 2026-10-21）（R24・R26・R28）。
%   (c) LLM との比較を、DB を渡して ID で判定する形（版 3、ゼロショット、4 モデル）に改めた。結果はまだない（本人の決定
%       2026-10-21: 旧版の Sonnet 4.5 の直接生成と等価判定の結果は本文と表から消す）（R27）。
%   (d) 主指標を Recall@20 にし、上位 K 件の指標を R-Precision とよぶ（旧 Recall@K）（R30）。別の出典の同じ式を正解と
%       みなした値も出す（R29）。自由度ゼロ停止の集合の指標（完全一致率・集合 F1・停止条件の充足率）は削除した。
%   (e) 検定をやめ、10 通りの分割のうち何通りで上回ったかを書く（旧 §5.5 Statistical Testing を削除）（R35）。
% 2026-10-21(2)（報告 f1076cd にそろえた）: 設定 A は 1,731 件。手法の設定値を hyperparameters とよぶ（§5.2 の題も）。
%   データ作成の LLM に、出典横断 113 件の説明文を書き直した Claude Opus 4.8 を加えた。
% ============================================================
\section{Evaluation Settings}
\label{sec:settings}

We evaluate the methods in two settings. Setting~A holds all 1,731 evaluation cases of the three
families. Setting~B holds the 1,000 synthesized cases alone, 100 cases for each number of
ground-truth equations from 1 to 10. The cases hold 3.99 ground-truth equations on average in
Setting~A and 5.50 in Setting~B. Neither setting uses the augmented cases
(Section~\ref{sec:cases}), and every result uses the stratified source-disjoint split of
Section~\ref{sec:split} with the same 10 random seeds.
% source: Dataset.tex 表 tab:families と同じ実測（synthesized 平均 5.50、設定 A 平均 3.99）

\section{Compared Methods and Hyperparameters}
\label{sec:compared}

We compare the proposed method with the classical information-retrieval baseline, two variants of
the baseline with stronger text scorers, and large language models (LLMs).

The classical baseline ranks the whole database by the Stage-1 score of Eq.~\eqref{eq:stage1},
without a learned model and without a candidate set. Its weights decide how well it ranks, so we
compare 12 weight pairs $(w_{\mathrm{t}}, w_{\mathrm{v}})$ from $(1.0, 0.0)$ to $(0.0, 1.0)$ and,
in each split, take the pair with the highest Recall@20 (Section~\ref{sec:metrics}) on the
validation cases to evaluate on the test cases. Setting~A chose $(0.4, 0.6)$ and $(0.3, 0.7)$ in
five splits each, and Setting~B chose $(0.3, 0.7)$ in five, $(0.2, 0.8)$ in three, and $(0.4, 0.6)$
in two; \ref{app:weights} lists all 12 pairs.
% source: experiments/kato1020/baselines.json（kato1020_baselines.py。検証用の Recall@20 で選択、同じなら MAP）;
%         main.json selected_weights（A tfidf: w40-60 5・w30-70 5、B tfidf: w30-70 5・w20-80 3・w40-60 2）

Two variants test whether stronger text evidence alone would close the gap. The first replaces the
TF-IDF cosine of Eq.~\eqref{eq:stage1} with BM25\cite{BM25} ($k_1 = 1.5$, $b = 0.75$, word
unigrams); the second replaces it with the cosine between dense embeddings of the description and
of each equation text from the sentence encoder E5\cite{E5} (\texttt{intfloat/e5-base-v2}: 12
layers, 768-dimensional embeddings, mean pooling), which was trained by contrastive learning for
retrieval. Because BM25 scores are unbounded and E5 cosines crowd into a narrow band, we rescale
each text score to $[0, 1]$ over the database for every query before mixing it with the
variable-overlap term. The validation cases choose the weights of both variants in the same way,
and both rank the whole database and train nothing.
% source: external_baselines.py（BM25 は自前実装: Lucene 形式 IDF ln(1+(N-n+0.5)/(n+0.5))、E5 は transformers で
%         平均プーリング＋L2 正規化、接頭辞 query:/passage:、最大 512 トークン）; kato1020_baselines.py SCORERS
%         （tfidf:none, bm25:minmax, e5:minmax。e5 のクエリは説明文だけ＝case_text(c)）

The LLMs receive the whole database and the query of one test case and return the identifiers of
up to 20 equations for the desired model, from the most to the least likely. They see no solved
example. Each equation of the database reaches them only as its identifier, its expression, and
its domain label, because adding the variable definitions and the context texts would exceed their
input limit of one million tokens. We withhold the sources, shuffle the order of the equations,
and renumber their identifiers: the proposed method does not read the sources, and the synthesized
ground truths prefer equations of one source, so the sources would give the LLMs alone a clue to
the answer. The models are Claude Haiku~5.5, Claude Sonnet~5.5, Claude Opus~5.5, and Claude
Fable~5.1, which differ in scale, all with the reasoning effort set to high
(\ref{app:prompts} gives the prompt). An LLM sees no training case, so its input for a
case does not depend on the split; we therefore ask it once per case and use the answer in every
split in which the case is a test case. We have prepared this comparison but have not yet run it.
% source: run_llm_db_baseline.py（版 3: fix-order・count3・pilot3・run3・evaluate3、commit 322fc73 / 81e206e）;
%         docs/appendix_materials/prompts/10_llm_db_baseline_v3_run_llm_db_baseline.txt。本実行は先生の返事と本人の承認待ち。
%         旧版（claude-sonnet-4-5 の直接生成、claude-opus-4-8 と claude-fable-5-1 の等価判定、各設定 50 件）の結果は
%         2026-10-21 の本人の決定で本文と表から消した（git 履歴の 1d8c26a 以前を参照）

The proposed method is the two-stage method of Chapter~\ref{Method}. We compared four first-stage
weight pairs, $(0.7, 0.3)$, $(0.5, 0.5)$, $(0.3, 0.7)$, and $(0.0, 1.0)$, by the Recall@20 of the
validation cases. With hard negatives, their validation Recall@20 ranged only from 0.980 to 0.983
in Setting~A and from 0.974 to 0.975 in Setting~B, a spread smaller than the variation across
splits (standard deviation about 0.015; \ref{app:weights}). We therefore use, in every split, the pair $(0.3, 0.7)$, whose
candidate set admits every ground-truth equation with the smallest $N_1$. In each split we set
$N_1$ to the smallest number of candidates that holds every ground-truth equation of the training
and validation cases, rounded up to the next thousand; under $(0.3, 0.7)$ this gives
$N_1 = 2{,}000$ in all 10 splits, whereas the other pairs need 2,000 to 6,000. Between random and
hard negatives, each split takes the one with the higher validation Recall@20, and hard negatives
won in 19 of the 20 splits of the two settings. The remaining hyperparameters --- a reference set of
$m = 5$ equations, the margin $\gamma = 0.1$, $n_{\mathrm{neg}} = 8$ incorrect candidates, and 15
epochs of AdamW with learning rate $10^{-3}$ and weight decay $10^{-4}$ --- are values we have used
since the beginning of this work and have not compared with others.
% source: experiments/kato1020/selected.json（--fix-weight w30-70。A: w30-70・hard・2,000 が 9 分割、w30-70・random・2,000 が
%         1 分割（7777）、B: w30-70・hard・2,000 が 10 分割。fixed: m 5, gamma 0.1, n_neg 8, epochs 15, hidden 64,
%         lr 1e-3, weight_decay 1e-4, batch 16）; selected_persplit.json grid（hard の検証用 Recall@20: A 0.980〜0.983、
%         B 0.974〜0.975）; n1_rule.json（(0.3, 0.7) は全分割で 2,000、ほかの重みは 2,000〜6,000）

\section{Evaluation Metrics}
\label{sec:metrics}

Every method returns identifiers of database equations: the retrieval methods rank them, and the
LLMs list the identifiers of the equations they were given. A returned equation is correct when its
identifier belongs to the ground truth $\mathcal{R}_q$ (Section~\ref{sec:formulation}); an
identifier absent from the database counts as incorrect, and a repeated identifier counts only at
its first occurrence. We judge no mathematical equivalence. The identifier-based judgment counts an
equation identical to a correct one but taken from another source as incorrect, although it would
serve the engineer equally well, and 386 of the 3,132 ground-truth equations (12.3\%) have such a
twin (Section~\ref{sec:cases}). We therefore also report the results with such twins counted as
correct; when several twins of one ground-truth equation appear, only the first counts.
% source: set_aware_reranker.py --equiv（ranks_equiv: latex_key＝gt_checks の正規化、別の出典の式だけ、1 本の正解式は 1 回）;
%         kato1020_baselines.equiv_from_pos（全 DB を並べる baseline に同じ規則）

For each test case we compute the following quantities and average them over the test cases of a
split; $\mathrm{Top}_n(q)$ denotes the top $n$ equations that a method returns.
\begin{itemize}
  \item Recall@20 $= |\mathcal{R}_q \cap \mathrm{Top}_{20}(q)| / K$, the primary metric: the
    fraction of the ground-truth equations among the 20 equations presented.
  \item Mean average precision (MAP): the mean over cases of the average precision
    \begin{equation}
      \mathrm{AP}(q) = \frac{1}{K} \sum_{e \in \mathcal{R}_q} P(e) ,
      \label{eq:ap}
    \end{equation}
    where $P(e)$ is the fraction of ground-truth equations among the equations ranked at or above
    $e$, and $P(e) = 0$ when $e$ is not in the ranking. MAP shows how closely the ground-truth
    equations crowd at the top. Because an LLM returns only 20 equations, the comparison with LLMs
    uses MAP@20, which sets $P(e) = 0$ for every ground-truth equation outside the top 20.
  \item R-Precision $= |\mathcal{R}_q \cap \mathrm{Top}_K(q)| / K$: the fraction of the
    ground-truth equations within the top $K$, where $K$ is the number of ground-truth
    equations\cite{Manning}.
  \item Candidate coverage $= |\mathcal{R}_q \cap C_q| / K$: the fraction of the ground-truth
    equations among the Stage-1 candidates.
\end{itemize}
% source: evaluate_multi_eq.py case_metrics（recall_at_20、R-Precision = recall_at_n_correct）; MAP は標準の定義
%         （kato1014・kato1020 の集計は保存した順位から計算する。並べた中にない正解式は P(e) = 0）; MAP@20 は上位 20 件で打ち切り
%         （2026-10-15 本人の追加依頼）。旧版の主指標 Recall@K は R-Precision とよぶ（加藤先生の注釈 R30）

We run no statistical tests. To compare two methods, we report the mean of their difference over
the 10 splits and the number of splits in which one beats the other.
% 2026-10-20 本人の決定（加藤先生の注釈 R35「検定自体、必要だとは思っていない」）。旧 §5.5 Statistical Testing
%   （Wilcoxon の符号順位検定、対応のある t 検定、Cohen の d_z）は削除した。

\section{Implementation and Reproducibility}
\label{sec:impl}

Every experiment runs as a deterministic script: the random seed fixes the split, the model
initialization, and the choice of negatives, so a rerun reproduces the reported numbers exactly.
The proposed method itself calls no language model at inference time. Language models appear only
in building the data and in the comparison with LLMs: Gemini extracted the equations, Gemini 2.5
Pro generated the reference sets and picked the single-source ground truths, the coding assistant
Claude wrote 33 cross-source combinations into a program, Claude Opus 4.8 rewrote the descriptions
of 113 cross-source cases, and Claude Sonnet 4.5 wrote the descriptions of the synthesized cases. The scripts specify their models, but no record keeps the
Gemini version used for the 9 earliest documents. \ref{app:hyper} lists every hyperparameter,
and the repository versions the scripts, the data, and the seeds.
% source: 抽出は Dataset.tex §4.1 と同じ（すべて Gemini、本人確認 2026-10-21）; 合成ケースの説明文は claude-sonnet-4-5-20250929
%         （dae_generation_report.json）
